\documentclass[12pt]{article}
\usepackage[utf8]{inputenc}
\usepackage[T1]{fontenc}
\usepackage{hyperref}
\title{Bolgarian Language Resources for K-12 Curriculum Development}
\author{Compiled by Claude Sonnet 5.5 (claude-sonnet-5-5) for Andr\'as Kornai}
\date{}
\begin{document}
\maketitle

\section{Introduction}
Bolgarian (ISO 639-3 xbo) is the Bulgar language, an extinct Oghur Turkic language. It is not the same as modern Bulgarian, which is Slavic. Bulgar died out in Danube Bulgaria around the 9th--10th century, as the Bulgar nobility became Slavicized. It persisted in the Volga region until about the 14th century. Chuvash is regarded as its living relative. Evidence is scarce: Volga Bulgar is known mainly from 13th--14th century tombstone inscriptions in Arabic script, and Danube Bulgar from a few loanwords in Old Bulgarian and from Greek inscriptions. The reference description said ``family: Turkic'', but the language has no living speakers. There is therefore no basis for a full curriculum in the language itself.

\section{Primary Educational Resources}
I found no primers, textbooks or school programmes for Bulgar in search results.

\section{Linguistic and Reference Materials}
\begin{itemize}
\item Wikipedia, ``Bulgar language'': overview of the Danube and Volga varieties and their classification. \url{https://en.wikipedia.org/wiki/Bulgar_language}
\item Verbix, Bulgar and Volga Bulgar pages. They are word lists and not authoritative. \url{https://docs.verbix.com/Languages/Bulgar} and \url{https://docs.verbix.com/Languages/BulgarVolga}
\item Wikipedia, ``Volga Bulgaria'': historical context. \url{https://wikipedia.com/wiki/Bulgar_Ulus}
\item Ayse Dietrich, ``Volga Bulgars language'' (Humanities Institute PDF): a general-audience overview. \url{https://humanitiesinstitute.org/__static/c0c8faf8456c0900d36c85ba3a3a91ff/volga-bulgars-language(2).pdf?dl=1}
\end{itemize}

\section{Religious and Community Resources}
I found no community or religious resources specific to Bulgar. Chuvash and Volga Tatar communities in Russia are historically connected to Volga Bulgaria, but I did not research their programmes here.

\section{Digital and Audio Resources}
I found no audio recordings or software. The inscriptions are the main primary source, but I saw no digital corpus.

\section{Curriculum Development Context}
Realistically, a K-12 unit would be history and heritage content: the Bulgars, Volga Bulgaria, the Bulgar calendar, loanwords, the Oghur and Chuvash connection, and how tombstone inscriptions are read. Teaching Bulgar as a spoken language is not possible with the evidence found. Chuvash materials would be a separate project under its own ISO code.
\end{document}